\documentclass{otengineering}
\newcommand{\breakprobe}[1]{%
    \par\medskip
    \textbf{Continuation probe #1.}
    This deliberately repeated paragraph makes the record long enough to
    require continuation onto another page. The frame and contents should
    remain complete, readable, and clear of the footer.
}
\begin{document}

\section{Generic box interface}

\begin{otbox}[Generic OT Box]{OTDark}
    This box directly tests the generic \texttt{otbox} interface. It should
    have a light-grey background, dark frame, bold title, sharp corners, and
    the established border thickness.
\end{otbox}

\section{Core engineering environments}

\begin{idea}
    The default title should be \textbf{Idea}, with the established blue
    frame.
\end{idea}

\begin{questionbox}
    The default title should be \textbf{Engineering Question}, with the
    established purple frame.
\end{questionbox}

\begin{decision}
    The default title should be \textbf{Design Decision}, with the
    established green frame.
\end{decision}

\begin{experiment}
    The default title should be \textbf{Experiment}, with the established
    orange frame.
\end{experiment}

\begin{failure}
    The default title should be \textbf{Failure / Problem}, with the
    established red frame.
\end{failure}

\begin{discovery}
    The default title should be \textbf{Discovery}, with the established
    purple frame.
\end{discovery}

\begin{lesson}
    The default title should be \textbf{Lesson Learned}, with the
    established green frame.
\end{lesson}

\begin{futurememo}
    The default title should be \textbf{Note to Future Self}, with the
    established orange frame.
\end{futurememo}

\section{Informal research and development environments}

\begin{thought}
    The default title should be \textbf{Random Thought}, with the
    established muted-grey frame.
\end{thought}

\begin{rabbithole}
    The default title should be \textbf{Rabbit Hole Tracker}, with the
    established orange frame.
\end{rabbithole}

\begin{researchqueue}
    The default title should be \textbf{Future Research Queue}, with the
    established blue frame.
\end{researchqueue}

\begin{wisdom}
    The default title should be \textbf{What I Wish I Knew Earlier}, with
    the established green frame.
\end{wisdom}

\section{Risk, assumption, and calculation environments}

\begin{risk}
    The default title should be \textbf{Risk Register Entry}, with the
    established red frame.
\end{risk}

\begin{assumption}
    The default title should be \textbf{Assumption Tracker}, with the
    established orange frame.
\end{assumption}

\begin{calculation}
    The default title should be \textbf{Engineering Calculation}, with the
    established blue frame.

    \calcfield{Given}{\(V=12\,\mathrm{V}\) and \(R=4\,\Omega\)}
    \calcfield{Required}{The current \(I\)}
    \calcfield{Calculation}{%
        \[
            I=\frac{V}{R}
             =\frac{12\,\mathrm{V}}{4\,\Omega}
             =3\,\mathrm{A}.
        \]
    }
    \calcfield{Result}{\(I=3\,\mathrm{A}\)}
\end{calculation}

\section{Custom titles}

\begin{idea}[Candidate Power-Supply Architecture]
    The optional argument should replace the default title while preserving
    the blue frame and all other established styling.
\end{idea}

\begin{failure}[Prototype Overheating]
    The optional argument should replace the default title while preserving
    the red frame and breakable behaviour.
\end{failure}

\begin{calculation}[Motor Starting-Current Estimate]
    The optional argument should replace the default calculation title.

    \calcfield{Nominal current}{\(I_{\mathrm{N}}=2.0\,\mathrm{A}\)}
    \calcfield{Starting-current factor}{\(k=5\)}
    \calcfield{Estimated starting current}{%
        \[
            I_{\mathrm{start}}
            =kI_{\mathrm{N}}
            =5(2.0\,\mathrm{A})
            =10\,\mathrm{A}.
        \]
    }
\end{calculation}

\clearpage
\section{Ordinary page-boundary behaviour}

The following space deliberately leaves too little room for the complete
box.

\rule{0pt}{0.92\textheight}

\begin{decision}[Boundary-Test Decision]
    This short box should move intact to the following page when there is
    insufficient space. Its title, background, contents, and frame should
    remain complete, with no overlap against the footer.
\end{decision}

\clearpage
\section{Breakable long-box behaviour}

The following space positions a long box near a page boundary.

\rule{0pt}{0.42\textheight}

\begin{experiment}[Long Breakable Experiment Record]
    \textbf{Objective.}
    Verify that the semantic wrapper inherits the generic box's breakable
    behaviour without losing its title or frame.

    \textbf{Initial condition.}
    The test begins near the bottom of a page, where the remaining vertical
    space cannot contain the complete record.

    \textbf{Observation 1.}
    The first portion should occupy the available space without overlapping
    the footer.

    \textbf{Observation 2.}
    The remaining material should continue on the following page with the
    established frame rendered correctly.

    \textbf{Observation 3.}
    Paragraph spacing and text alignment should remain readable on both
    pages.

    \textbf{Observation 4.}
    Mathematical content should remain aligned:
    \[
        P=VI,
        \qquad
        E=Pt,
        \qquad
        \eta=\frac{P_{\mathrm{out}}}{P_{\mathrm{in}}}.
    \]

    \textbf{Observation 5.}
    The page break should not clip any line, duplicate content, or create
    unexpected blank space.

    \textbf{Conclusion.}\breakprobe{1}
\breakprobe{2}
\breakprobe{3}
\breakprobe{4}
\breakprobe{5}
\breakprobe{6}
\breakprobe{7}
\breakprobe{8}
\breakprobe{9}
\breakprobe{10}
\breakprobe{11}
\breakprobe{12}
    Successful continuation across the page boundary confirms the current
    breakable behaviour of both the semantic wrapper and the underlying
    \texttt{otbox} interface.
\end{experiment}

\end{document}